\documentclass{beamer}

% Theme selection
\usetheme{Madrid}
\usecolortheme{default}

% Packages
\usepackage[utf8]{inputenc}
\usepackage[T1]{fontenc}
\usepackage{amsmath}
\usepackage{amssymb}
\usepackage{graphicx}
\usepackage{braket} % For Dirac notation

% Title information
\title[Konstruksi Hamiltonian QBGT]{Konstruksi Hamiltonian pada Quantum Bayesian Game Theory (QBGT) untuk Optimasi Portofolio}
\subtitle{Derivasi Matematis dari Teori Permainan ke Operator VQE}
\date{\today}

\begin{document}

% Slide 0: Title
\begin{frame}
    \titlepage
\end{frame}

% Slide 1: Definisi Masalah Bayesian
\begin{frame}{Definisi Masalah Bayesian}
    Pasar saham dapat dipandang sebagai permainan strategi antara Investor melawan "Nature" dengan informasi tidak lengkap (\textit{incomplete information}).
    
    \begin{itemize}
        \item \textbf{Himpunan Tipe Nature:}
        Alam memiliki dua kemungkinan kondisi:
        $$ \Theta = \{\theta_{bull}, \theta_{bear}\} $$
        
        \item \textbf{Probabilitas Prior (Belief $q$):}
        Merepresentasikan keyakinan investor terhadap kondisi pasar.
        \begin{itemize}
            \item \textbf{Statis (QBGT):} Bertindak sebagai bobot dalam konstruksi \textit{Grand Hamiltonian}.
            \item \textbf{Dinamis:} Sebagaimana evolusi belief $q_t$ pada \textit{Chain Store Paradox}, dimana investor memperbarui keyakinan berdasarkan sinyal pasar (Bayesian Update).
        \end{itemize}
        $$ P(\theta_{bull}) = q, \quad P(\theta_{bear}) = 1-q $$
        
        \item \textbf{Matriks Payoff:}
        Keuntungan (Return) bergantung pada kondisi pasar yang terjadi:
        \begin{itemize}
            \item Kondisi Bullish: Matriks $M_{bull}$ (High Return)
            \item Kondisi Bearish: Matriks $M_{bear}$ (Loss/Low Return)
        \end{itemize}
    \end{itemize}
\end{frame}

% Slide 2: Skema Kuantum EWL
\begin{frame}{Skema Kuantum EWL (Eisert-Wilkens-Lewenstein)}
    Protokol permainan kuantum standar memetakan strategi klasik ke dalam operasi unitary di ruang Hilbert.
    
    \begin{enumerate}
        \item \textbf{State Awal:} Qubit dimulai dari keadaan dasar.
        $$ |\psi_0\rangle = |00\dots0\rangle $$
        
        \item \textbf{Entanglement:} Operator $\hat{J}$ menciptakan korelasi kuantum antar pemain/aset.
        $$ |\psi_{ent}\rangle = \hat{J} |\psi_0\rangle $$
        
        \item \textbf{Strategi Unitary:} Pemain (Investor) menerapkan strategi $\hat{U}$.
        $$ |\psi_{strat}\rangle = (\hat{U}_A \otimes \hat{U}_B) |\psi_{ent}\rangle $$
        
        \item \textbf{State Akhir:} Setelah disentangling operation $\hat{J}^\dagger$.
        $$ |\psi_f\rangle = \hat{J}^\dagger (\hat{U}_A \otimes \hat{U}_B) \hat{J} |00\dots0\rangle $$
    \end{enumerate}
\end{frame}

% Slide 3: Integrasi Bayesian ke dalam EWL (The QBGT Model)
\begin{frame}{Integrasi Bayesian ke dalam EWL (The QBGT Model)}
    Dalam kerangka QBGT, pemain tidak menghadapi satu matriks payoff tunggal, melainkan superposisi dari beberapa kemungkinan payoff berdasarkan tipe Nature.
    
    \vspace{0.3cm}
    \textbf{Ekspektasi Payoff Kuantum:}
    Nilai harap keuntungan ($\Pi$) adalah jumlahan terbobot dari nilai ekspektasi operator pengukuran payoff ($\hat{\$}$) untuk setiap kondisi pasar:
    
    \begin{block}{Persamaan Ekspektasi Payoff}
    $$ \Pi(\hat{U}) = q \cdot \langle \psi_f | \hat{\$}_{bull} | \psi_f \rangle + (1-q) \cdot \langle \psi_f | \hat{\$}_{bear} | \psi_f \rangle $$
    \end{block}
    
    Dimana:
    \begin{itemize}
        \item $\hat{\$}_{bull/bear}$ adalah operator Hermisian yang merepresentasikan matriks payoff $M_{bull/bear}$.
        \item $q$ adalah kepercayaan investor terhadap kondisi pasar Bullish.
    \end{itemize}
\end{frame}

% Slide 4: Konstruksi Hamiltonian (Bagian 1 - Pembobotan)
\begin{frame}{Konstruksi Hamiltonian (Bagian 1 - Pembobotan)}
    Untuk menggunakan algoritma VQE (Variational Quantum Eigensolver), masalah \textit{Payoff Maximization} harus diubah menjadi \textit{Energy Minimization}.
    $$ H \propto -\Pi $$
    
    Parameter portofolio (Return $\mu$ dan Risiko $\sigma$) dikombinasikan secara linear \textit{sebelum} dikonstruksi menjadi Hamiltonian sirkuit:
    
    \begin{columns}
        \begin{column}{0.48\textwidth}
            \begin{block}{Effective Return}
            $$ \mu_{eff} = q\mu_{bull} + (1-q)\mu_{bear} $$
            \end{block}
        \end{column}
        \begin{column}{0.48\textwidth}
            \begin{block}{Effective Risk (Covariance)}
            $$ \sigma_{eff} = q\sigma_{bull} + (1-q)\sigma_{bear} $$
            \end{block}
        \end{column}
    \end{columns}
    
    \vspace{0.3cm}
    Ini memungkinkan VQE mencari satu solusi robust yang sudah memperhitungkan ketidakpastian kedua skenario pasar.
\end{frame}

% Slide 5 (Baru): Formulasi Fungsi Biaya (Cost Function)
\begin{frame}{Formulasi Fungsi Biaya Markowitz \& Penalti}
    Sebelum dipetakan ke kuantum, kita definisikan fungsi objektif klasik yang menggabungkan tujuan investasi dan kendala.
    
    $$ O(x) = \underbrace{w \sum_{i,j} \sigma_{ij}^{eff} x_i x_j}_{\text{Minimalkan Risiko}} - \underbrace{(1-w) \sum_i \mu_i^{eff} x_i}_{\text{Maksimalkan Return}} + \underbrace{P \left(\sum_i x_i - B\right)^2}_{\text{Kendala Budget}} $$

    \begin{itemize}
        \item \textbf{$x_i \in \{0,1\}$:} Variabel keputusan (1 jika aset $i$ dipilih, 0 jika tidak).
        \item \textbf{$w$ (Risk Aversion):} Parameter preferensi investor ($0 \le w \le 1$). 
            \begin{itemize}
                \item $w=1$: Hanya peduli risiko (Min Variance).
                \item $w=0$: Hanya mengejar profit (Max Return).
            \end{itemize}
        \item \textbf{$\sigma_{ij}^{eff}$:} Matriks kovarians efektif (gabungan Bull/Bear).
        \item \textbf{$\mu_i^{eff}$:} Vektor return efektif.
        \item \textbf{$B$:} Target jumlah aset dalam portofolio (Cardinality Constraint).
        \item \textbf{$P$:} Penalti energi yang sangat besar untuk memaksa solusi mematuhi kendala $B$.

    \end{itemize}
\end{frame}

% Slide 5: Konstruksi Hamiltonian (Bagian 2 - Mapping ke Ising)
\begin{frame}{Konstruksi Hamiltonian (Bagian 2 - Mapping ke Ising)}
    Optimasi Markowitz menggunakan variabel biner $x_i \in \{0, 1\}$. Komputer kuantum menggunakan operator spin Pauli-Z $Z_i \in \{+1, -1\}$.
    
    \textbf{Transformasi Variabel:}
    $$ x_i \rightarrow \frac{1-Z_i}{2} $$
    
    Substitusi ke dalam fungsi objektif $O(x)$:
    $$ O(x(\vec{Z})) = w \sum_{ij} \sigma_{ij}^{eff} \left(\frac{1-Z_i}{2}\right) \left(\frac{1-Z_j}{2}\right) - \dots + P(\dots)^2 $$
    
    Ekspansi aljabar menghasilkan bentuk Ising yang siap untuk VQE:
    $$ H \sim \sum_i h_i Z_i + \sum_{ij} J_{ij} Z_i Z_j $$
\end{frame}

% Slide 6: Hasil Akhir - Hamiltonian Sistem
\begin{frame}{Hasil Akhir - Hamiltonian Sistem}
    Setelah ekspansi aljabar, kita mendapatkan Hamiltonian Ising final yang siap dijalankan pada VQE:
    
    \begin{alertblock}{Ising Hamiltonian}
    $$ H_{final} = \sum_i h_i Z_i + \sum_{i<j} J_{ij} Z_i Z_j + C $$
    \end{alertblock}
    
    \begin{itemize}
        \item \textbf{Longitudinal Field ($h_i$):} Bias lokal pada setiap qubit. Merepresentasikan return aset individu + efek penalti linear + bobot Bayesian.
        \item \textbf{Coupling ($J_{ij}$):} Interaksi antar qubit. Merepresentasikan risiko kovarians antar aset + efek penalti kuadratik.
        \item \textbf{Konstanta ($C$):} Pergeseran energi global (bisa diabaikan dalam optimasi).
    \end{itemize}
    
    \textit{Ground State} dari Hamiltonian ini adalah strategi portofolio yang paling optimal dan robust terhadap ketidakpastian pasar.
\end{frame}

% Slide 7: Sintesis - The Giant Operator
\begin{frame}{Sintesis: Operator Hamiltonian Raksasa}
    Dalam konteks QBGT, Hamiltonian final ini sebenarnya adalah "Operator Raksasa" yang menggabungkan seluruh semesta kemungkinan.
    
    \begin{block}{The Grand Hamiltonian}
    $$ \hat{H}_{QBG} = \sum_{\theta \in \Theta} P(\theta) \cdot \hat{H}_{\theta} $$
    \end{block}
    
    Untuk kasus dua kondisi pasar (Bull/Bear), operator ini menjadi superposisi linear dari Hamiltonian spesifik per-kondisi:
    
    $$ \hat{H}_{QBG} = q \cdot \hat{H}_{Bull} + (1-q) \cdot \hat{H}_{Bear} $$
    
    \begin{itemize}
        \item $\hat{H}_{Bull}$: Operator energi yang "menghukum" strategi defensif saat pasar naik.
        \item $\hat{H}_{Bear}$: Operator energi yang "menghukum" strategi agresif saat pasar turun.
        \item $\hat{H}_{QBG}$: Operator kompromi yang mencari \textit{Ground State} penengah (Robust Strategy).
    \end{itemize}
    
    Inilah inti dari penggabungan Game Theory dan Fisika Kuantum: \textbf{Menemukan keseimbangan energi di antara realitas yang bertentangan.}
\end{frame}

\end{document}
